\section{Introduction}
\label{sec:introduction}

%\todo[inline]{TODO(Toma): pregled podrucja}
Unmanned aerial vehicles (UAVs) have become a popular solution for performing autonomous inspection of various structures, such as windturbines. The attractive aspect of this approach is the ability to perform non destructive testing (NDT) without putting people at risk. This has led to the development of UAV-based inspection techniques in a range of tasks. Traditionally, wind turbine inspection involved either the dangerous task of personnel climbing the windturbine or in some cases the expensive task of unmounting the rotor blades for inspection.\\
In this work, a method is proposed for performing autonomous visual inspection of the rotor blades using a UAV while the turbine is rotating. Previous work on visual inspection of stationary wind turbines involves the need to halt normal wind turbine operation, making it inconvenient and expensive. Performing inspection on a rotating wind turbine, on the other hand, circumvents this problem. 


Several previous studies have explored the use of UAVs, mainly quadrotors, for the inspection of stationary wind turbines. For example, vision based inspection systems of wind turbines are presented in articles such as~\cite{weibin_gu, stokkeland2015, schafer2016, ref:Car2020}.

In~\cite{weibin_gu} a vision based system for UAV wind turbine inspection is designed and the results presented both in simulation and Hardware-In-The-Loop testing. The system uses a LiDAR sensor and a camera for the vision pipeline integrating the YOLOv3 network and a customized Hough transform algorithm.\\
Furthermore, in~\cite{stokkeland2015} a machine vision module for the navigation of an UAV during the inspection of stationary wind turbines is presented. It implements the estimation of the relative position and distance between the UAV and the wind turbine, and the position of the blades. The system utilized the Hough transform for detection and the Kalman filter for tracking. Experiments show the accuracy and robustness of the solution.\\
Once again,~\cite{schafer2016} studies autonomous inspection flights at a wind turbine, focusing on the generation of an a-priori 3D map of the wind turbine, and path planning and collision avoidance algorithms. The system relies on a GPS and a 2D LiDAR sensor to collect point clouds which are used for the relative localization process.\\
Moreover, in~\cite{ref:Car2020} the team on whose work this study is based on presents a semi-autonomous wind turbine blade inspection system with a LiDAR-equipped UAV. The process performs successful wind turbine blade inspections with minimal operator involvement and results in blade images and a wind turbine 3D model. The method is tested and validated in a real life setting.\\
Lastly, in~\cite{drewry} a review of NDT techniques for wind turbines is given, giving more insight into the actual inspection techniques rather than the manner in which they are carried out (manually or automated).


When talking about the use of UAVs for the inspection of rotating wind turbines, however, the research is very limited. This is a more challenging task due to the dynamic nature of the moving turbine and the need to maintain a safe distance of the UAV from the rotating blades at all times. Nevertheless, some recent studies have shown promising results in this area.\\
For example, in~\cite{doi:10.1177/0309524X211027814} a novel method is proposed where an UAV is used to perform strain monitoring on the wind turbine rotor based on a optical digital image correlation technique. Validation of the method is performed using strain gauges proving it robust.\\
Furthermore, in~\cite{KHADKA2020106446} a digital image correlation system installed on a drone is used to obtain dynamic characteristics of the blades with high precision while they rotate, enabling on-site robust measurements without the need to stop the operation of the wind turbine.\\

Despite these advances, the following technical challenges need to be addressed to perform UAV inspection of rotating wind turbines. These are: developing robust algorithms to detect and track the blades, ensure safe and reliable operation of the UAV in the challenging and dynamic environment, and keeping track of information acquired dynamically for each blade.\\

In this article, we present the work done on the problem of rotating wind turbine inspection using a UAV. The implemented system contains the method for finding the relative position of the rotating wind turbine relative to the position of the UAV, relying heavily on the LiDAR sensor system, an inertial measurement unit, and a GPS. Moreover, once the relative position is found, we developed a method for performing non destructive visual analysis of the rotating rotor blades by planning a suitable trajectory and triggering a camera with a gimbal camera system. Having implemented the aforementioned techniques, the function of the system was experimentally tested on data collected on a real world wind farm.\\

The article is divided into sections as follows. Section \ref{sec:introduction} provides an introduction into the studied matter, including an overview of available literature focusing on similar problems. The 3D model generation method used to obtain the point cloud of the rotating wind turbine and the model matching algorithm used to find the generated model in the data stream coming from the LiDAR sensors is presented in Section \ref{sec:model_matching}. In Section \ref{sec:camera_triggering} the trajectory generation algorithm to perform the inspection the rotor blades of a rotating wind turbine and the setup allowing for autonomous camera triggering is presented. The following Section \ref{sec:experimental_results} presents the results of experiments performed using the presented approach on data collected in real world, followed by a short discussion of the results. Funally, Section \ref{sec:conclusion} contains a conclusion of the study, with a short recap of the work and possible avenues for future work.

%\todo[inline]{TODO(lmark): Odrediti doprinose}
%\textbf{Doprinosi (prijedlog):}
%\begin{enumerate}
%    \item Pronalazak rotirajuce vjetroelektrane u mapi (opis algoritma)
%    \item Automatsko trigeriranje kamere + planiranje trajektorije za snimanje rotirajuce vjetroelektrane
%    \item Eksperimentalna provjera sustava
%\end{enumerate}
%
%\textbf{Opis instroductiona:}
%\begin{enumerate}
%    \item Opis problema
%    \item Pregled podrucja (LiDAR, kamere, ostali senzori)
%    \item Pregled podrucja (stacionarne, rotirajuce, itd.)
%    \item Istaknuti u cemo smo drugaciji, tj. 2-3 doprinosa tokom clanka
%    \item Predstaviti ostatak clanka - Kasnije
%\end{enumerate}